\phantomsection
\section*{专题——溶液中的微粒浓度比较}
\addcontentsline{toc}{section}{专题——溶液中的微粒浓度比较}
\label{sect:专题——溶液中的微粒浓度比较}

        \begin{xmp}
        [专题——溶液中的微粒浓度比较例 1]
        {}
        []
        []
        在 \SI{0.1}{\mole\per\litre} 的 \cemh{NH3.H2O} 溶液中，下列关系正确的是（\makebox[2em]{A}）
        \begin{itemize}
            \item[A.] \(c(\cemh{NH3.H2O}) > c(\cemh{OH^-}) > c(\cemh{NH4+}) > c(\cemh{H+})\)
            \item[B.] \(c(\cemh{NH4+}) > c(\cemh{NH3.H2O}) > c(\cemh{OH^-}) > c(\cemh{H+})\)
            \item[C.] \(c(\cemh{NH3.H2O}) > c(\cemh{NH4+}) = c(\cemh{OH^-}) > c(\cemh{H+})\)
            \item[D.] \(c(\cemh{NH3.H2O}) > c(\cemh{NH4+}) > c(\cemh{H+}) > c(\cemh{OH^-})\)
        \end{itemize}
    \end{xmp}

    \begin{xmp}
        [专题——溶液中的微粒浓度比较例 2]
        {}
        []
        []
        \SI{0.1}{\mole\per\litre} 的醋酸钠溶液中

        （1）含有\uline{\makebox[3em]{7}}种微粒（分子和离子）；

        （2）分别是：\uline{\makebox[24em]{\cemh{CH3COO^-}、\cemh{Na+}、\cemh{CH3COOH}、\cemh{H2O}、\cemh{OH^-}、\cemh{H+}}}；

        （3）按浓度从大到小的顺序排列上述微粒：（除 \cemh{H2O}）
        
        \uline{\makebox[28em]{\(c(\cemh{Na+}) > c(\cemh{CH3COO^-}) > c(\cemh{CH3COOH}) > c(\cemh{OH^-}) > c(\cemh{H+})\)}}；

        （4）下列关于该溶液中微粒浓度的关系是否正确：
        \begin{itemize}
            \item[A.] \(c(\cemh{Na+}) + c(\cemh{H+}) = c(\cemh{CH3COO^-}) + c(\cemh{OH^-})\) \dotfill\footnote{电荷守恒}（\makebox[1.5em]{\hcmark}）
            \item[B.] \(c(\cemh{Na+}) = c(\cemh{CH3COO^-}) + c(\cemh{CH3COOH})\)             \dotfill\footnote{物料守恒}（\makebox[1.5em]{\hcmark}）
            \item[C.] \(c(\cemh{CH3COO^-}) + c(\cemh{CH3COOH}) = \SI{0.1}{\mole\per\litre}\) \dotfill（\makebox[1.5em]{\hcmark}）
            \item[D.] \(c(\cemh{H+}) + c(\cemh{CH3COOH}) = c(\cemh{OH^-})\)                  \dotfill\footnote{质子守恒}（\makebox[1.5em]{\hcmark}）
        \end{itemize}
    \end{xmp}

    \begin{xmp}
        [专题——溶液中的微粒浓度比较例 3]
        {}
        []
        []
        \SI{0.1}{\mole\per\litre} 的 \cemh{NaHS} 溶液中（已知：\cemh{H2S} 的 \(K_{\text{a}1}=\num{1.1E-7}, K_{\text{a}2}=\num{1.3E-13}\)）

        （1）含有\uline{\makebox[3em]{7}}种微粒（分子和离子）；

        （2）分别是：\uline{\makebox[24em]{\cemh{H2S}、\cemh{HS^-}、\cemh{S^2-}、\cemh{Na+}、\cemh{H2O}、\cemh{H+}、\cemh{OH^-}}}；

        （3）按浓度从大到小的顺序排列上述微粒：（除 \cemh{H2O}）

        \uline{\makebox[28em]{\(c(\cemh{Na+}) > c(\cemh{HS^-}) > c(\cemh{OH^-}) > c(\cemh{H2S}) > c(\cemh{H+}) > c(\cemh{S^2-})\)}}；

        （4）下列关于该溶液中微粒浓度的关系是否正确：
        \begin{itemize}
            \item[A.] \(c(\cemh{Na+}) + c(\cemh{H+}) = c(\cemh{HS^-}) + c(\cemh{OH^-})\)     \dotfill\footnote{\(+2c(\cemh{S^2-})\)}（\makebox[1.5em]{\hxmark}）
            \item[B.] \(c(\cemh{Na+}) + c(\cemh{H+}) = c(\cemh{HS^-}) + c(\cemh{S^2-}) + c(\cemh{OH^-})\)  \dotfill（\makebox[1.5em]{\hxmark}）
            \item[C.] \(c(\cemh{Na+}) = c(\cemh{HS^-}) + c(\cemh{S^2-}) = \SI{0.1}{\mole\per\litre}\)      \dotfill\footnote{\(-c(\cemh{H2S})\)}（\makebox[1.5em]{\hxmark}）
            \item[D.] \(c(\cemh{H+}) = c(\cemh{HS^-}) + c(\cemh{OH^-})\)   \dotfill（\makebox[1.5em]{\hxmark}）
            \item[E.] \(c(\cemh{OH^-}) = c(\cemh{H2S}) + c(\cemh{H+})\)    \dotfill（\makebox[1.5em]{\hxmark}）
        \end{itemize}
    \end{xmp}

    \begin{xmp}
        [专题——溶液中的微粒浓度比较例 4]
        {}
        []
        []
        在 \SI{0.1}{\mole\per\litre} 的 \cemh{Na2CO3} 溶液中，下列关于微粒浓度的关系是否正确：
        \begin{itemize}
            \item[A.] \(c(\cemh{Na+}) + c(\cemh{H+}) = c(\cemh{CO3^2-}) + c(\cemh{HCO3^-}) + c(\cemh{OH^-})\) \dotfill（\makebox[1.5em]{\hxmark}）
            \item[B.] \(c(\cemh{Na+}) > c(\cemh{CO3^2-}) > c(\cemh{OH^-}) > c(\cemh{HCO3^-}) > c(\cemh{H+})\) \dotfill（\makebox[1.5em]{\hcmark}）
            \item[C.] \(c(\cemh{Na+}) = 2[c(\cemh{HCO3^-}) + c(\cemh{CO3^2-}) + c(\cemh{H2CO3})]\) \dotfill（\makebox[1.5em]{\hcmark}）
            \item[D.] \(c(\cemh{OH^-}) = c(\cemh{H+}) + c(\cemh{HCO3^-}) + c(\cemh{H2CO3})\) \dotfill（\makebox[1.5em]{\hxmark}）
            \item[E.] \(c(\cemh{Na+}) + c(\cemh{OH^-}) = c(\cemh{H+}) + 2c(\cemh{CO3^2-}) + 3c(\cemh{HCO3^-}) + 4c(\cemh{H2CO3})\) \\
            \null\dotfill（\makebox[1.5em]{\hcmark}）
        \end{itemize}
    \end{xmp}

    \begin{xmp}
        [专题——溶液中的微粒浓度比较例 5]
        {}
        []
        []
        在物质的量浓度均为 \SI{0.01}{\mole\per\litre} 的 \cemh{HCN} 和 \cemh{NaCN} 混合溶液中，测得 \(c(\cemh{Na+}) > c(\cemh{CN^-})\)，则下列关系正确的是（\makebox[2em]{BD}）（不定项）
        \begin{itemize}
            \item[A.] \(c(\cemh{H+}) > c(\cemh{OH^-})\)
            \item[B.] \(c(\cemh{H+}) < c(\cemh{OH^-})\)
            \item[C.] \(c(\cemh{Na+}) > c(\cemh{HCN})\)
            \item[D.] \(c(\cemh{HCN}) + c(\cemh{CN^-}) = \SI{0.02}{\mole\per\litre}\)
        \end{itemize}
    \end{xmp}

    \begin{xmp}
        [专题——溶液中的微粒浓度比较例 6]
        {}
        []
        []
        把 \SI{0.02}{\mole\per\litre} \cemh{CH3COOH} 溶液和 \SI{0.01}{\mole\per\litre} \cemh{NaOH} 溶液等体积混和，则混合液中微粒浓度关系正确的为（\makebox[2em]{AD}）（不定项）
        \begin{itemize}
            \item[A.] \(c(\cemh{CH3COO^-}) > c(\cemh{Na+}) > c(\cemh{H+}) > c(\cemh{OH^-})\)
            \item[B.] \(c(\cemh{Na+}) > c(\cemh{CH3COO^-}) > c(\cemh{H+}) > c(\cemh{OH^-})\)
            \item[C.] \(c(\cemh{CH3COOH}) > c(\cemh{CH3COO^-})\)
            \item[D.] \(c(\cemh{CH3COOH}) + c(\cemh{CH3COO^-}) = \SI{0.01}{\mole\per\litre}\)
            \item[E.] \(2c(\cemh{H+}) = c(\cemh{CH3COO^-}) - c(\cemh{CH3COOH})\)
        \end{itemize}
    \end{xmp}

    \phantomsection
    \subsection*{归纳：如何比较溶液中的微粒浓度}
    \addcontentsline{toc}{subsection}{归纳：如何比较溶液中的微粒浓度}
    \label{subsec:归纳：如何比较溶液中的微粒浓度}

        \begin{enumerate}
            \item 若发生反应，则先计算\uuline{混合}后各溶质的浓度，再拆成微粒浓度\textcolor{red}{（此时暂不考虑电离和水解）}；若不发生反应则跳过此步。
            \item 考虑溶液中\textcolor{blue}{\uline{所有}}平衡对微粒浓度的影响（弱电解质电离、盐类水解以及水的电离）
            
            其中： \Circled{1} 先考虑主过程，再考虑次过程；\textcolor{red}{（主过程 \(>\) 次过程）}
            
            \phantom{其中：}\Circled{2} 弱电解质电离、盐类水解都是微弱的，水的电离更弱
            
            \phantom{其中：\Circled{2}} \textcolor{red}{（一般：电离 \(>\) 水解）}

            \item 守恒关系：
            
            电荷守恒：一边全是正离子，另一边全是负离子，\textcolor{red}{记得要乘以电荷数}
            
            物料守恒：一边来源于正离子，另一边来源于负离子，
            
            \phantom{物料守恒：}\textcolor{red}{记得微粒要找全}
            
            质子守恒：由前两个守恒叠加得到

            \item 若全部都是一价离子，由电荷守恒结合溶液酸碱性，可推导离子浓度的大小关系
        \end{enumerate}

    \begin{xmp}
        [专题——溶液中的微粒浓度比较特例]
        {特例}
        []
        []
        \SI{25}{\celsius} 时，\(K_{\text{a}1}(\cemh{H2S}) = \num{1.1E-7}, K_{\text{a}2}(\cemh{H2S}) = \num{1.3E-13}\)。此温度下，\SI{0.01}{\mole\per\litre} \cemh{Na2S} 溶液的有关微粒浓度关系正确的是\uline{\makebox[3em]{BC}}。（不定项）

        \begin{itemize}
            \item[A.] \(c(\cemh{Na+}) + c(\cemh{H+}) < c(\cemh{OH^-}) + c(\cemh{HS^-}) + c(\cemh{S^2-})\)
            \item[B.] \(c(\cemh{Na+}) > c(\cemh{OH^-}) > c(\cemh{HS^-}) > c(\cemh{S^2-})\)
            \item[C.] \(c(\cemh{Na+}) + c(\cemh{H+}) = c(\cemh{OH^-}) + c(\cemh{HS^-}) + 2c(\cemh{S^2-})\)
            \item[D.] \(c(\cemh{Na+}) = 2c(\cemh{S^2-}) + 2c(\cemh{HS^-})\)
        \end{itemize}

        \vspace{0.5\baselineskip}

        \begin{tabular}{r@{\quad}c@{}c@{}c@{}c@{}c@{}c@{}c}
            & \(\mathbf{\cemh{S^2-}} \) & \(\mathbf{\cemh{+}}\)  & \(\mathbf{\cemh{H2O}}\)  & \(\mathbf{\cemh{<=>}} \) & \(\mathbf{\cemh{HS^-}} \) & \(\mathbf{\cemh{+}} \) & \(\mathbf{\cemh{OH^-}} \) \\
            平衡浓度（\si{\mole\per\litre}） & \(\mathcolor{blue}{0.01-x}\) &&&& \(\mathcolor{blue}{x}\) && \(\mathcolor{blue}{x}\) \\
        \end{tabular}

        \[
            K_{\text{h}1} = \frac{K_\text{W}}{K_{\text{a}2}} = \frac{\num{1E-14}}{\num{1.3E-13}} = 0.0769
        \]
        \[
            K_\text{h} = \frac{x^2}{0.01 - x} \quad \text{解得} x= 0.00896
        \]
        \cemh{S^2-} 的水解程度接近 \SI{90}{\percent}，大部分都水解了

        \(\therefore c(\cemh{OH^-}) > c(\cemh{HS^-}) > c(\cemh{S^2-})\)
    \end{xmp}
